\section{Gauge and circuit holonomy}

A local change of sheet labels is a function
\[
g:V\longrightarrow\Ztwo
\]
acting by
\[
G_g(i,\epsilon)=(i,\epsilon+g_i).
\]
Conjugating (T_v) by (G_g) produces another voltage vector (v') with
\[
v_i'=v_i+g_i+g_{i+1}.
\]
This is fiber-preserving gauge equivalence.

\begin{definition}
The circuit holonomy of (v) is
\[
\hol(v)=\sum_{i=0}^{n-1}v_i\in\Ztwo.
\]
\end{definition}

\begin{proposition}
Circuit holonomy is invariant under gauge transformation.
\end{proposition}

\begin{proof}
Summing the transformed voltages gives
\[
\sum_i v_i'
=\sum_i v_i+\sum_i g_i+\sum_i g_{i+1}
=\sum_i v_i,
\]
because the last two sums agree and cancel in (Ztwo).
\end{proof}

\begin{proposition}
Two voltage assignments are gauge equivalent if and only if they have the
same circuit holonomy.
\end{proposition}

\begin{proof}
Necessity was just proved. For sufficiency, suppose (v) and (w) have the
same sum. Set (g_0=0) and recursively define
\[
g_{i+1}=g_i+v_i+w_i
\]
for (0\leq i<n-1). Then (w_i=v_i+g_i+g_{i+1}) for the first (n-1)
edges. Equality of the total sums supplies the same identity on the final
edge. Hence (v) and (w) are gauge equivalent.
\end{proof}

The holonomy is therefore not merely an invariant. It is the complete gauge
classifier.

